%% kapa: exam=1 round=36 year=2025 period=1 form=A subject=economics subject_ko=경제학원론 seq=17 number=57
%% source: Q-Net 감정평가사 기출문제 (cst00309) artlSeq=5231525 file=raw/1008/5231525/2025년 제36회 감정평가사 1차시험 1교시 문제지 원본.pdf
%% conversion: kordoc 4.24.51 -> tools/gen1.py (auto); review=proofread-sonnet; flags=none
\begin{question}{57}
정보의 비대칭성에 관한 설명으로 옳은 것은?
\begin{choices}
  \item 본인-대리인 문제는 감추어진 유형(hidden type)으로 인해 발생하는 문제에 해당한다.
  \item 정부가 가입을 강제하는 건강보험과 같은 단체보험은 도덕적 해이 문제를 완화할 수 있다.
  \item 품질보증은 역선택 문제를 해결하는 방안 중 하나이다.
  \item 보험시장에서 기초공제(initial deduction)는 역선택 문제를 완화할 수 있다.
  \item 애컬로프(Akerlof)의 중고차시장 모형은 도덕적 해이 문제를 설명한 것이다.
\end{choices}
\end{question}
\answer{3}
